\documentclass[11pt]{article}

\usepackage{fullpage}
\usepackage{amsmath}
\usepackage{epsfig}
\usepackage{graphics}
\usepackage{caption}
\usepackage{circuitikz}

\oddsidemargin -0.15in
\evensidemargin -0.15in
\headheight 0in
\headsep 0in
%\topmargin -0.25in
\textheight 9.3in
\textwidth 6.8in

\begin{document}


\begin{center}
\section*{
Franklin \& Marshall College \hfill PHY 223\\
\bigskip
{\Large \it Speed of light}}
\end{center}

The speed of light in free space is an important and intriguing constant of
nature. Whether the light comes from a laser on a desktop or from a distant
galaxy that is hurtling away at fantastic speeds, the measurement of the speed
of light will yield the same constant value. In more precise terminology, the
speed of light is independent of the relative velocities of the light source and
the observer.

As Einstein first presented in his special theory of relativity, the speed of
light is critically important in some surprising ways:

\begin{enumerate}
 \item The speed of light establishes an upper limit to the speed that may be 
 imparted to any material object.
 
 \item Objects moving near the speed of light follow a set of physical laws 
 drastically different, not only from Newton’s Laws, but from the basic 
 assumptions of human intuition.
\end{enumerate}

It is not surprising that a great deal of time and effort has been invested in
measuring the speed of light. Some of the most accurate measurements were made
by Albert Michelson between 1926 and 1929. Michelson measured the speed of light
in air to be 2.99712x10$^8$ m/s. From this result he deduced the speed of light
in free space to be 2.99796x10$^8$ m/s. Today, the speed of light in free space
is, by definition, $(2.997924562 \pm 0.000000011)x10^8$ m/s exactly.

Our approach to the measurement of the speed of light will be direct and
intuitive, and simpler then Michelson's experiment. We will use a diode laser as
our light source. We will send pulses of light across the lab, measure how long
they take to come back, and deduce the speed of light from the knowledge that
light travels at constant speed.

\vspace{5mm}

\noindent {\bf Apparatus}

\begin{itemize}
  
\item Diode Laser — The diode laser emits an intense, narrowly- focused beam of
light at a wavelength of approximately 650nm. In this experiment, the diode
laser is powered by a function generator, which modulates the light intensity at
about 4 MHz. The laser is equipped with adjustment screws for precisely aiming
the light at the mirror.

\item Function Generator — Provides power to the diode laser. The laser is
turned on by a DC offset voltage from the function generator and modulated by a
superimposed sine wave having frequency of about 4 MHz.

\item Concave Mirror — The concave surface of the mirror helps to focus the
light as it is reflected. The mirror is also equipped with adjustment screws for
aiming the light back to the receiver.

\item Converging Lens — The lens is used to focus returning light onto the
active element of the light receiver. The lens has a focal length of 12.7cm.

\item Light Receiver — The receiver is designed for receiving audio and video
signals transmitted via modulated light. Since the light receiver is sensitive
to very high-frequency modulation, it is ideally suited to this experiment.
There are two sensitive elements on the receiver; in this experiment, you will
use only the one labeled “Video” (which is sensitive to high frequency signals).
(The “Audio” receiver, which is sensitive to lower frequency signals, will not
be used in this experiment.)

\item Oscilloscope — The oscilloscope is used to view both the modulating signal
and the returned signal simultaneously. The phase difference between these two
will be converted to a time delay.

\item Other Equipment — Pasco OS-9172 Alignment bench, cables, tape measure.
\end{itemize}

\vspace{5mm}

\noindent {\bf Laser light safety notes}

We will use a relatively safe, low power laser. Nevertheless, take the following
precautions:

\begin{enumerate}

\item Never look directly into the laser beam, either directly, or as it is
reflected from a mirror or other surface.

\item Keep the optical path at floor level and keep your eyes well above floor
level at all times.

\item Only people involved with the experiment should be nearby.

\end{enumerate}

\vspace{5mm}

\noindent {\bf Experimental setup}

\vspace{3mm}

\begin{minipage}{\linewidth} \centering
\includegraphics[scale=1.5]{./Experimental_Setup.png}
  \captionof{figure}{ Experimental set-up showing the laser, concave mirror,
lens, and receiver. The laser, lens, and receiver are mounted on the laser
alignment bench. The mirror is shown mounted on a stand, but in our arrangement
both the bench and the mirror will be placed directly on the laboratory floor.}
\end{minipage}

\vspace{3mm}

\begin{enumerate}
  
\item Place the laser, lens, and receiver on the alignment bench as shown in the
figure above.
  
\item Place the bench on the floor, making sure that you have 12-15 meters of
available space in front of it.

\item Connect your function generator to the diode's input power and to channel
1 of the oscilloscope. This will be the signal you trigger from.

\item Connect the video output of the light receiver to channel 2 of the
oscilloscope.

\item Set the function generator for a sine wave with a frequency as close to 4
MHz as you can make it, and press the DC Offset button to turn the DC offset
function on. Turn the Output and DC Offset knobs completely counterclockwise to
the zero position.

\item Power on the function generator and turn up the DC Offset knob until you
see laser light.
  
\item Align the laser so that it propagates along a line on the floor and that
its height is constant throughout. It will make alignment easier afterwards.
Once this alignment is done,

\item Place the mirror approximately 1 m in front of the laser and use it to
reflect the beam (through the lens) on the light receiver sensor. Make sure you
go through the center of the lens. It will make the procedure simpler.

\item Check the strength of the signal on the oscilloscope. You want to see a
sine wave with decent amplitude (signal-to-noise). You might have to play with
the function generator some more (amplitude and DC offset) to get a better
signal. Set the oscilloscope to average 128 cycles.
  
\end{enumerate}

\vspace{5mm}

\noindent {\bf Procedure}

\begin{itemize}
\item Using a tape measure, record the total path length of the light pulses.

\item Using the cursors on the oscilloscope, measure the time delay between the
two traces.

\item Also using the cursors, also record your uncertainty in the time delay.

\item Move the mirror back about 60 cm and repeat the measurement.
\end{itemize}


\vspace{5mm}

\noindent {\bf Analysis}

Make a graph of the time delay as a function of round trip distance. Find the
slope of this graph and its uncertainty. Convert this slope (and uncertainty)
into a value for the speed of light (and uncertainty).


% \bibliographystyle{plain}

% \begin{thebibliography}{10}

% \bibitem {tag} L. J. van der Pauw, ``A Method of Measuring Specific
% Resistivity and Hall Effect of Discs of Arbitrary Shape,'' Philips Research
% Reports 13, 1 (1958).

% \end{thebibliography}


\end{document}

%%% Local Variables:
%%% mode: latex
%%% TeX-master: t
%%% End:
